\section{Related Works}
\label{sec:related_works}
\textbf{Pooling architectures for permutation invariant mappings} \hspace{1pt}
Pooling architectures for sets have been used in various problems such as 3D shape recognition~\citep{BaoguangShi2015, Su2015},
discovering causality~\citep{Lopez-Paz2016},
learning the statistics of a set~\citep{Edwards2017},
few-shot image classification~\citep{Snell2017},
and conditional regression and classification~\citep{Garnelo2018}.
\citet{Zaheer2017} discuss the structure in general and provides a partial proof of the universality of the pooling architecture,
and \citet{Wagstaff2019} further discuss the limitation of pooling architectures. \citet{Bloem-Reddy2019} provides a link between
probabilistic exchangeability and pooling architectures.

\textbf{Attention-based approaches for sets} \hspace{1pt}
Several recent works have highlighted the competency of attention mechanisms in modeling sets.
\citet{Vinyals2016a} pool elements in a set by a weighted average with weights computed using an attention mechanism.
\citet{Yang2018} propose AttSets for multi-view 3D reconstruction, where dot-product attention is applied to compute the weights used to pool the encoded features via weighted sums.
Similarly, \citet{Ilse2018} use attention-based weighted sum-pooling for multiple instance learning.
Compared to these approaches, ours use multihead attention in aggregation, and more importantly, we propose to apply self-attention after pooling to model correlation among multiple outputs.
PMA with $k=1$ seed vector and single-head attention roughly corresponds to these previous approaches.
Although not permutation invariant, \citet{Mishra2018} has attention as one of its core components to meta-learn to solve various tasks using sequences of inputs.
\citet{Kim2019} proposed attention-based conditional regression, where self-attention is applied to the query sets.

\textbf{Modeling interactions between elements in sets} \hspace{1pt}
An important reason to use the Transformer is to explicitly model higher-order interactions among the elements in a set.
\citet{Santoro2017} propose the relational network, a simple architecture that sum-pools all pairwise interactions of elements in a given set, but not higher-order interactions.
Similarly to our work, \citet{Ma2018} use the Transformer to model interactions between the objects in a video.
They use mean-pooling to obtain aggregated features which they fed into an LSTM.

\textbf{Inducing point methods} \hspace{1pt}
The idea of letting trainable vectors $I$ directly interact with data points is loosely based on the inducing point methods used in sparse Gaussian processes~\citep{Snelson2005} and the Nystr\"{o}m method for matrix decomposition~\citep{Fowlkes2004}.
$m$ trainable inducing points can also be seen as $m$ independent memory cells accessed with an attention mechanism.
The differential neural dictionary~\citep{Pritzel2017} stores previous experience as key-value pairs and uses this to process queries.
One can view the ISAB is the inversion of this idea, where queries $I$ are stored and the input features are used as key-value pairs.

%\ak{we might want to cite https://arxiv.org/abs/1805.00613 as a different approach to the problem: instead of considering permutation invariance, we could just treat permutation as a latent variable}
